\clearpage
\fancyhead[L]{\small INHALT.}
\pdfbookmark[0]{Inhalt der Ausgabe von 1874}{contents1874}
\SourceLeaf{660}
\newcommand{\GaussContentsEntry}[1]{%
  \par\nointerlineskip
  \vbox{%
    \hsize=\linewidth\relax
    \raggedright\parskip=0pt
    \noindent\strut#1\strut\par
  }%
  \smallskip
}
\newcommand{\GaussContentsLeaders}[1]{%
  \unskip\nobreak\hskip .33em\relax#1%
}
\newcommand{\GaussContentsReference}[1]{%
  \unskip\nobreak\hskip .7em plus 1fill\relax\mbox{#1}%
}
\begingroup
\raggedright
\begin{center}\textbf{INHALT.}\par\medskip
\textbf{GAUSS WERKE BAND VI. ASTRONOMISCHE ABHANDLUNGEN.}\end{center}
\begin{center}\small Die Seitenzahlen beziehen sich auf die Ausgabe von 1874.\end{center}
\par\medskip
\Needspace{5\baselineskip}\noindent\textit{Abhandlungen.}\par\smallskip\nopagebreak[4]
\GaussContentsEntry{\noindent\hangindent=1em Disquisitio de elementis ellipticis Palladis ex oppositionibus annorum 1803, 1804, 1805, 1807, 1808, 1809\GaussContentsLeaders{\mbox{.\ .\ .\ .\ .\ .\ .\ .\ .\ .\ .\ .\ .\ .\ .\ .}}\GaussContentsReference{\mbox{1810 Nov. 25}\quad\mbox{Seite 1}}}
\GaussContentsEntry{\noindent\hangindent=1em Observationes cometae secundi a. 1813 in observatorio Gottingensi factae, adjectis nonnullis adnotationibus circa calculum orbitarum parabolicarum\GaussContentsReference{\mbox{1813 Sept. 10}\quad\mbox{--- 25}}}
\GaussContentsEntry{\noindent\hangindent=1em Methodus peculiaris elevationem poli determinandi\GaussContentsLeaders{\mbox{.\ .\ .\ .\ .\ .\ .\ .\ .}}\GaussContentsReference{\mbox{1808}\quad\mbox{--- 37}}}
\par\medskip
\Needspace{5\baselineskip}\noindent\textit{Anzeigen eigener Schriften.}\par\smallskip\nopagebreak[4]
\GaussContentsEntry{\noindent\hangindent=1em Theoria motus corporum coelestium in sectionibus conicis Solem ambientium\GaussContentsReference{\mbox{1809 Juni 17}\quad\mbox{--- 53}}}
\GaussContentsEntry{\noindent\hangindent=1em Disquisitio de elementis ellipticis Palladis ex oppositionibus annorum 1803, 1804, 1805, 1807, 1808, 1809\GaussContentsLeaders{\mbox{.\ .\ .\ .\ .\ .\ .\ .\ .\ .\ .\ .\ .\ .\ .\ .}}\GaussContentsReference{\mbox{1810 Dez. 10}\quad\mbox{--- 61}}}
\GaussContentsEntry{\noindent\hangindent=1em Observationes cometae secundi a. 1813 in observatorio Gottingensi factae, adjectis nonnullis adnotationibus circa calculum orbitarum parabolicarum\GaussContentsReference{\mbox{1814 Jan. 3}\quad\mbox{--- 65}}}
\GaussContentsEntry{\noindent\hangindent=1em Methodus peculiaris elevationem poli determinandi\GaussContentsLeaders{\mbox{.\ .\ .\ .\ .\ .\ .\ .\ .}}\GaussContentsReference{\mbox{1808 Dez. 5}\quad\mbox{--- 68}}}
\par\medskip
\Needspace{5\baselineskip}\noindent\textit{Verschiedene Aufsätze über Astronomie.}\par\smallskip\nopagebreak[4]
\GaussContentsEntry{\noindent\hangindent=1em Berechnung des Osterfestes\GaussContentsLeaders{\mbox{.\ .\ .\ .\ .\ .\ .\ .\ .\ .\ .\ .\ .\ .\ .\ .\ .\ .}}\GaussContentsReference{\mbox{1800 Aug.}\quad\mbox{--- 73}}}
\GaussContentsEntry{\noindent\hangindent=1em Berechnung des Pessachfestes\GaussContentsLeaders{\mbox{.\ .\ .\ .\ .\ .\ .\ .\ .\ .\ .\ .\ .\ .}}\GaussContentsReference{\mbox{1802 Mai}\quad\mbox{--- 80}}}
\GaussContentsEntry{\noindent\hangindent=1em Noch etwas über die Bestimmung des Osterfestes\GaussContentsLeaders{\mbox{.\ .\ .\ .\ .\ .\ .\ .\ .\ .}}\GaussContentsReference{\mbox{1807 Sept. 12}\quad\mbox{--- 82}}}
\GaussContentsEntry{\noindent\hangindent=1em Vorschriften, wie man aus der geozentrischen Länge und Breite eines Himmelskörpers, der Position seines Knotens, der Neigung seiner Bahn, der Länge der Sonne und ihrem Abstand von der Erde die heliozentrische Länge des Himmelskörpers in seiner Bahn, seinen wahren Abstand von der Sonne und seinen wahren Abstand von der Erde ableitet\GaussContentsLeaders{\mbox{.\ .\ .\ .\ .\ .\ .\ .\ .\ .\ .}}\GaussContentsReference{\mbox{1802 Juni}\quad\mbox{--- 87}}}
\GaussContentsEntry{\noindent\hangindent=1em Einige Bemerkungen zur Vereinfachung der Berechnung der geozentrischen Positionen der Planeten\GaussContentsLeaders{\mbox{.\ .\ .\ .\ .\ .\ .\ .\ .\ .\ .\ .\ .\ .\ .\ .\ .\ .\ .}}\GaussContentsReference{\mbox{1804 Mai}\quad\mbox{--- 94}}}
\GaussContentsEntry{\noindent\hangindent=1em Über die Grenzen der geozentrischen Positionen der Planeten\GaussContentsLeaders{\mbox{.\ .\ .\ .\ .\ .}}\GaussContentsReference{\mbox{1804 Aug.}\quad\mbox{--- 106}}}
\GaussContentsEntry{\noindent\hangindent=1em Der Zodiakus der Juno\GaussContentsLeaders{\mbox{.\ .\ .\ .\ .\ .\ .\ .\ .\ .\ .\ .\ .\ .\ .\ .\ .\ .\ .}}\GaussContentsReference{\mbox{1805 März}\quad\mbox{--- 119}}}
\endgroup
